\chapter{Sets}

\section{Why This Matters}
Sets are incredibly powerful for two main reasons:
\begin{enumerate}
    \item \textbf{Deduplication:} The fastest way to remove duplicates from a list is to convert it into a set.
    \item \textbf{Membership Testing:} Because sets use a data structure called a "hash table" under the hood, checking if an item exists in a set (e.g., \texttt{if "apple" in my\_set:}) is almost instantly fast, even if the set has millions of items. Doing this with a list is incredibly slow.
\end{enumerate}

\section{Explanation}
A set in Python is an unordered, mutable collection of \textbf{unique} items. You can think of it exactly like a mathematical set. Because they are unordered, \textbf{they do not support indexing or slicing} (you cannot do \texttt{my\_set[0]}).

\subsection{1. Syntax and Core Methods}
\begin{lstlisting}
unique_numbers = {1, 2, 3, 4, 4, 4, 5}
print(unique_numbers)  # {1, 2, 3, 4, 5} -> Duplicates removed!

# Creating an empty set
empty_set = set() # MUST use set(), because {} creates a dict!
\end{lstlisting}

\begin{itemize}
    \item \texttt{.add(item)}: Adds an item to the set.
    \item \texttt{.remove(item)}: Removes an item (throws an error if it doesn't exist).
    \item \texttt{.discard(item)}: Removes an item (does NOTHING if it doesn't exist - safer!).
\end{itemize}

\subsection{2. Mathematical Operations}
Sets shine when you need to compare two groups of data.
\begin{lstlisting}
group_a = {"apple", "banana", "cherry"}
group_b = {"cherry", "date", "elderberry"}

print(group_a | group_b) # Union: {'apple', 'banana', 'cherry', 'date', 'elderberry'}
print(group_a & group_b) # Intersection: {'cherry'}
print(group_a - group_b) # Difference: {'apple', 'banana'}
\end{lstlisting}

\subsection{3. The Hashable Requirement}
Sets can only store \textbf{immutable} (hashable) items. This means you can put strings, integers, and tuples inside a set, but you CANNOT put a list or a dictionary inside a set!
\begin{lstlisting}
valid_set = {1, "hello", (10, 20)} # Works perfectly
# invalid_set = {1, "hello", [10, 20]}  <--- TypeError!
\end{lstlisting}

\mlconnection{In Natural Language Processing (NLP), you often need to find the "vocabulary" of a document---the unique words used. If a document has 10,000 words, you simply run \texttt{vocab = set(words\_list)} to instantly get the unique words. Sets are also used to quickly identify overlapping features between datasets.}

\section{Solutions to Exercises}

\textbf{Exercise 1: Deduplication}
\begin{lstlisting}
emails = ["a@a.com", "b@b.com", "a@a.com", "c@c.com", "b@b.com"]
unique_emails = set(emails)
print(len(unique_emails)) # 3
\end{lstlisting}

\vspace{1em}
\textbf{Exercise 2: Intersection}
\begin{lstlisting}
user_likes = {"The Matrix", "Inception", "Interstellar"}
friend_likes = {"Interstellar", "The Prestige", "The Matrix"}
print(user_likes & friend_likes)
\end{lstlisting}

\vspace{1em}
\textbf{Exercise 3: The Empty Set Trap}
\begin{lstlisting}
unique_ips = set()
unique_ips.add("192.168.1.1")
\end{lstlisting}
\textit{Solution: \texttt{\{\}} evaluates to an empty dictionary, which does not have an \texttt{.add()} method. You must use \texttt{set()}.}
